\documentclass[11pt,a4paper]{article}

\usepackage[T1]{fontenc}
\usepackage{lmodern}
\usepackage[margin=2.4cm]{geometry}
\usepackage{enumitem}
\usepackage{booktabs}
\usepackage{tabularx}
\usepackage[hidelinks]{hyperref}
\usepackage{xurl}

\setlength{\parindent}{0pt}
\setlength{\parskip}{5pt}
\setlist[itemize]{leftmargin=1.6em,itemsep=2pt,topsep=3pt}

\begin{document}

\begin{center}
{\Large\bfseries DATASET DESCRIPTION}
\end{center}

This document describes the two public road-extraction datasets used to train
and evaluate DBR-Net: the Massachusetts Roads Dataset and the DeepGlobe Road
Extraction Dataset. Both datasets formulate road extraction as binary semantic
segmentation, but they differ in spatial resolution, geographical coverage,
scene composition, and annotation characteristics.

\section{Massachusetts Roads Dataset}

\textbf{Dataset source:} A public aerial-image dataset introduced by Volodymyr
Mnih in 2013 and distributed through the University of Toronto. The copy used
in this study was obtained from the Kaggle mirror listed in Section~\ref{sec:data-availability}.

\subsection{Overview of the Dataset}

The Massachusetts Roads Dataset is a widely used benchmark for extracting road
networks from high-resolution aerial imagery. It covers urban, suburban, and
rural areas across Massachusetts, USA. Each input image is paired with a binary
target map in which road pixels form the foreground class and all remaining
pixels form the background class. The large native images and the mixture of
dense street networks, major roads, and narrow local roads make the dataset
suitable for evaluating both regional road coverage and fine-structure recovery.

\subsection{Dataset Composition}

\begin{itemize}
  \item \textbf{Total size:} 1,171 aerial image--mask pairs.
  \item \textbf{Official training split:} 1,108 image--mask pairs.
  \item \textbf{Official validation split:} 14 image--mask pairs.
  \item \textbf{Official test split:} 49 image--mask pairs.
  \item \textbf{Image size:} $1500\times1500$ pixels.
  \item \textbf{Spatial resolution:} approximately 1~m/pixel.
  \item \textbf{Geographical coverage:} more than 2,600~km$^2$.
  \item \textbf{Data type:} three-channel optical aerial imagery with binary
        road target maps generated from vector road data.
  \item \textbf{Task:} pixel-level binary semantic segmentation of roads and
        background.
\end{itemize}

\subsection{Image and Annotation Characteristics}

\begin{itemize}
  \item \textbf{Scene diversity:} The images include dense residential areas,
        suburban neighbourhoods, highways, intersections, vegetation, water
        bodies, and sparsely developed regions.
  \item \textbf{Road morphology:} Road networks contain long linear segments,
        curved streets, junctions, loops, and narrow branches at multiple
        orientations and scales.
  \item \textbf{Occlusion and ambiguity:} Trees, buildings, shadows, and vehicles
        can obscure road surfaces, while rivers, roofs, and other elongated
        structures may resemble roads locally.
  \item \textbf{Class imbalance:} Road pixels occupy a substantially smaller
        image area than background pixels.
  \item \textbf{Native-image challenge:} The $1500\times1500$ resolution requires
        a model to preserve fine details while retaining sufficient contextual
        information over a large field of view.
\end{itemize}

\subsection{Use in the Present Study}

\begin{itemize}
  \item \textbf{Benchmark role:} Massachusetts Roads serves as the first
        benchmark for training and evaluating DBR-Net and the comparison models.
  \item \textbf{Training input:} Training uses random $1024\times1024$ crops from
        the native images.
  \item \textbf{Data augmentation:} Random horizontal and vertical flips,
        rotations by $90^{\circ}$, $180^{\circ}$, or $270^{\circ}$, and random
        scaling with a factor sampled from $[0.75,1.0]$ are applied.
  \item \textbf{Centerline supervision:} Auxiliary centerline targets are
        generated from the binary masks by tensor-based morphological thinning
        before downsampling; they are not separate annotations supplied by the
        dataset.
  \item \textbf{Evaluation:} Predictions are binarised at a threshold of 0.5.
        Precision, recall, F1-score, and road-class intersection over union (IoU)
        are computed from pixel-level confusion statistics.
\end{itemize}

\subsection{Dataset Strengths and Limitations}

\begin{itemize}
  \item \textbf{Strengths:} The dataset provides public pixel-level targets,
        large native images, varied road layouts, and a standard split that
        supports comparison with previous road-extraction studies.
  \item \textbf{Limitations:} Its geographical coverage is restricted to one US
        state, and the rasterized road labels may not represent the exact visible
        width of every road surface. Results on this dataset alone therefore do
        not establish generalisation to other countries, sensors, or annotation
        conventions.
\end{itemize}

\section{DeepGlobe Road Extraction Dataset}

\textbf{Dataset source:} The road-extraction track of the DeepGlobe 2018
Satellite Image Understanding Challenge. The copy used in this study was
obtained from the Kaggle mirror listed in Section~\ref{sec:data-availability}.

\subsection{Overview of the Dataset}

The DeepGlobe Road Extraction Dataset was developed for large-scale road
segmentation from satellite imagery. It contains geographically diverse scenes
from Thailand, Indonesia, and India and includes urban, suburban, rural, and
agricultural environments. Unlike centerline-only annotations, its binary masks
label the visible road area at the pixel level. The variation in road material,
density, surrounding land cover, and illumination makes DeepGlobe a challenging
benchmark for testing model robustness across heterogeneous scenes.

\subsection{Dataset Composition}

\begin{itemize}
  \item \textbf{Total size:} 8,570 RGB satellite images.
  \item \textbf{Original challenge split:} 6,226 training images, 1,243
        validation images, and 1,101 test images.
  \item \textbf{Publicly labelled subset:} 6,226 training image--mask pairs in
        the commonly distributed release; the original validation and test
        labels are not included in that release.
  \item \textbf{Image size:} $1024\times1024$ pixels.
  \item \textbf{Spatial resolution:} approximately 0.5~m/pixel.
  \item \textbf{Geographical coverage:} approximately 2,220~km$^2$.
  \item \textbf{Data type:} three-channel RGB satellite imagery with binary
        pixel-level road masks.
  \item \textbf{Class distribution:} Road pixels account for approximately
        4.5\% of the original training split, illustrating the strong foreground
        imbalance of the task.
\end{itemize}

\subsection{Image and Annotation Characteristics}

\begin{itemize}
  \item \textbf{Geographical diversity:} Samples cover multiple countries and
        include urban, suburban, rural, agricultural, forested, and coastal
        environments.
  \item \textbf{Road diversity:} The dataset contains paved and unpaved roads,
        highways, local streets, intersections, and road networks with widely
        varying density and width.
  \item \textbf{Pixel-level road areas:} The masks delineate road surfaces rather
        than only their centerlines, enabling direct road-area segmentation.
  \item \textbf{Challenging appearance:} Illumination changes, vegetation,
        shadows, low-contrast unpaved roads, and visually similar non-road
        structures increase ambiguity.
  \item \textbf{Foreground imbalance:} Most pixels belong to the background,
        making foreground-sensitive metrics such as road IoU and F1-score more
        informative than overall pixel accuracy alone.
\end{itemize}

\subsection{Use in the Present Study}

\begin{itemize}
  \item \textbf{Benchmark role:} DeepGlobe serves as the second benchmark for
        evaluating DBR-Net under broader geographical and scene variation.
  \item \textbf{Reproducible labelled split:} The 6,226 labelled pairs are
        shuffled with seed 3407. Exactly 5,000 samples are used for training and
        the remaining 1,226 samples form the full test holdout. The first 300
        samples of this holdout are also used for validation and remain within
        the reported 1,226-image test set.
  \item \textbf{Model-selection dependency:} Because the 300-image validation
        subset is contained in the test holdout, the DeepGlobe test score is not
        fully independent of model selection. This dependency must be considered
        when interpreting the result.
  \item \textbf{Training input:} The native image size is
        $1024\times1024$ pixels, matching the training crop size used in the main
        experiments.
  \item \textbf{Training and evaluation:} The same augmentation operations,
        optimisation policy, decision threshold, and pixel-level metrics are
        applied to DBR-Net and all comparison models.
\end{itemize}

\subsection{Dataset Strengths and Limitations}

\begin{itemize}
  \item \textbf{Strengths:} DeepGlobe offers large-scale, geographically varied
        satellite imagery and pixel-level road-area annotations covering diverse
        road surfaces and scene types.
  \item \textbf{Limitations:} Public labels are concentrated in the original
        training portion, which requires a custom split for reproducible local
        evaluation. The present study's overlapping validation and test protocol
        also introduces a model-selection dependency.
\end{itemize}

\section{Comparative Summary}

\begin{table}[h]
\centering
\small
\begin{tabularx}{\textwidth}{@{}lXX@{}}
\toprule
\textbf{Property} & \textbf{Massachusetts Roads} & \textbf{DeepGlobe Roads} \\
\midrule
Image modality & RGB aerial imagery & RGB satellite imagery \\
Native size & $1500\times1500$ pixels & $1024\times1024$ pixels \\
Spatial resolution & Approximately 1~m/pixel & Approximately 0.5~m/pixel \\
Total images & 1,171 & 8,570 \\
Geographical scope & Massachusetts, USA & Thailand, Indonesia, and India \\
Road annotation & Binary rasterized road target maps & Binary road-surface masks \\
Primary challenge & Large scenes and fine road structures & Cross-region and
appearance diversity \\
\bottomrule
\end{tabularx}
\end{table}

Together, the datasets provide complementary evaluation conditions. The
Massachusetts Roads Dataset tests extraction from large aerial tiles within a
consistent regional domain, whereas DeepGlobe tests robustness to broader
geographical, environmental, and road-surface variation.

\section{Data Availability}
\label{sec:data-availability}

The dataset copies used in this study are publicly accessible through Kaggle:

\begin{itemize}
  \item Massachusetts Roads Dataset:\newline
        \url{https://www.kaggle.com/datasets/balraj98/massachusetts-roads-dataset}
  \item DeepGlobe Road Extraction Dataset:\newline
        \url{https://www.kaggle.com/datasets/balraj98/deepglobe-road-extraction-dataset}
\end{itemize}

The original Massachusetts Roads Dataset and its citation information are
available from the University of Toronto at
\url{https://www.cs.utoronto.ca/~vmnih/data/}. The DeepGlobe dataset is described
in I. Demir et al., ``DeepGlobe 2018: A Challenge to Parse the Earth Through
Satellite Images,'' \textit{Proceedings of the IEEE/CVF Conference on Computer
Vision and Pattern Recognition Workshops}, 2018, pp.~172--181,
\url{https://openaccess.thecvf.com/content_cvpr_2018_workshops/html/Demir_DeepGlobe_2018_A_CVPR_2018_paper.html}.

\end{document}
